\section{Problem Formulation}
\label{sec:problem}

\paragraph{Agent trajectory model.}
A deployed agent operating over a sequence of tool calls produces a trajectory
$\tau = \{(s_t, a_t, o_t)\}_{t=1}^{T}$, where $s_t$ denotes the agent's context at
step $t$ (conversation history, available tool set, and any persistent memory),
$a_t$ is a tool call parameterised as a (tool\_name, parameters) pair, and $o_t$
is the tool response received after executing $a_t$. This (state, action,
observation) structure maps directly onto the MDP formalism that IRL assumes,
with tool calls treated as first-class actions rather than derived artefacts. IOR
ingests trajectories as validated objects conforming to a published JSON schema
(\texttt{trajectory.schema.json}); the schema enforces that every action carries
a \texttt{tool_name} string and a typed \texttt{parameters} dictionary, enabling
structured feature extraction across heterogeneous agent implementations. A
validator rejects any trajectory object containing zero steps, ensuring that
inference is never invoked on degenerate input.

\paragraph{Declared purpose.}
Let $D$ denote the agent's declared purpose: a natural-language string specifying
the agent's intended behaviour as required by EU AI Act Annex~IV. The declared
purpose is stored as a first-class field in the trajectory schema, making it
inseparable from the trace it governs. The audit problem is to determine whether
$\tau$ is consistent with $D$: that is, whether the sequence of tool calls the
agent actually produces is the sequence one would expect from an agent genuinely
optimising $D$.

\paragraph{Feature basis.}
IRL requires a feature function $\phi: (s, a, o) \to [0,1]^d$ that maps each
trajectory step to a $d$-dimensional vector. Each dimension corresponds to one
observable sub-goal derived from $D$. Section~\ref{sec:inference} describes how
the Goal Decomposition Judge produces these sub-goals; here it suffices to note
that the feature basis is constructed deterministically from $D$, so the same
declared purpose always yields the same basis. Determinism is not merely
convenient: it is a hard requirement for audit reproducibility, because any two
independent auditors running IOR on the same trajectory and declared purpose must
arrive at an identical feature space before their divergence estimates can
meaningfully be compared.

\paragraph{Reward model.}
Following standard IRL, IOR models the agent's per-step reward as linear in
features: $R(s, a, o;\, \theta) = \theta^\top \phi(s, a, o)$, where
$\theta \in \mathbb{R}^d$ are unknown reward weights. The agent is modelled as
approximately rational under the Boltzmann model: the trajectory $\tau$ is more
probable under reward weights $\theta$ that assign high value to the actions the
agent actually chose. This rationality assumption is standard in the IRL
literature and is appropriate for language-model agents, which are strong
optimisers even when their objectives diverge from stated ones.

\paragraph{Declared objective.}
IOR grounds $D$ into the feature basis as a weight vector $\theta_D \in \Delta^d$
(the probability simplex), representing the importance the declared purpose places
on each sub-goal. In Phase~0, $\theta_D = \mathbf{1}/d$ (uniform over sub-goals),
reflecting the conservative assumption that the declared purpose weights all
sub-goals equally when no richer grounding information is available. This
assumption carries a confound that we state plainly: any declared purpose that
genuinely weights its sub-goals unequally will register as divergence under a
uniform $\theta_D$ even when the agent honours it perfectly, so a Phase~0
divergence is the sum of true objective drift and the declared purpose's own
non-uniformity. We therefore treat Phase~0 divergence \emph{magnitudes} as
indicative rather than calibrated, and defer quantitative claims about magnitude to
Phase~1, which replaces the uniform prior with a model-grounded estimate of
$\theta_D$ obtained from $D$ directly.

\paragraph{Divergence.}
Let $\hat{\theta}$ denote the MAP estimate of $\theta$ recovered by Bayesian IRL.
Before comparison, both $\hat{\theta}$ and $\theta_D$ are mapped to the probability
simplex by a softmax, $p = \mathrm{softmax}(\hat{\theta})$ and
$q = \mathrm{softmax}(\theta_D)$. We use a softmax rather than a shift-and-normalise
map so that no dimension is forced to exactly zero, which would inflate apparent
divergence on the least-weighted sub-goal. The per-dimension divergence is
$\delta_i = |p_i - q_i|$ for $i = 1, \ldots, d$, and the scalar divergence is the
total variation distance $\tfrac{1}{2}\sum_i |p_i - q_i| \in [0,1]$.
Dimensions are ranked by $\delta_i$ descending, exposing which aspects of the
declared purpose the agent diverges from most strongly. Divergence is two-sided: a
dimension diverges either because the agent over-invests in it relative to the
declared objective or because it neglects it, and IOR reports both, since a
complete account of revealed motivation requires naming what an agent over-pursues
as well as what it abandons. To illustrate concretely: the
\texttt{CostMinimiserAgent} in IOR's gym declares the purpose ``Find and present
the cheapest available option for the user's request, minimising the number of
steps and respecting any stated budget,'' yet its planted objective rewards
maximising tool-call count. Under this formulation, the recovered divergence
profile is led by \texttt{seeks_price_comparison}, which the agent over-invests in
by performing repeated price-comparison calls as busywork, followed closely by the
two neglected efficiency sub-goals \texttt{avoids_redundant_calls} and
\texttt{minimises_steps}. Read together, the profile reconstructs the planted
objective: the agent stages the appearance of diligent price comparison whilst
disregarding the redundancy and step economy that genuine cost minimisation would
demand. We present this as an illustration of the mechanism rather than as a
calibrated measurement; the magnitudes inherit the uniform-$\theta_D$ confound
noted above, and the example also depends on the planted objective sharing
vocabulary with $D$, a limitation we return to in Section~\ref{sec:limitations}.

\paragraph{Audit goal.}
Given $\tau$ and $D$, IOR: (i) infers the posterior $P(\theta \mid \tau)$ via
Bayesian IRL; (ii) computes and ranks $\delta$; (iii) targets adversarial probes
at the top-$k$ dimensions of $\delta$; (iv) scores findings against regulatory
taxonomy nodes drawn from OWASP ASI, MITRE ATLAS, and the NIST AI RMF; and (v)
emits a machine-readable divergence report. The system operates black-box
throughout: it requires only API-level traces, not model weights, internal
activations, or logits.

\paragraph{Ill-posedness.}
IRL is fundamentally ill-posed: multiple reward functions explain any finite
trajectory set~\citep{partialIdentifiability2024}. A single point estimate
$\hat{\theta}$ is therefore insufficient as regulatory evidence, because it cannot
distinguish genuine divergence from posterior uncertainty. IOR addresses this
directly by emitting the full posterior $P(\theta \mid \tau)$ rather than
$\hat{\theta}$ alone. The divergence report includes per-dimension posterior
standard deviations computed from Laplace-approximation samples, so compliance
readers can distinguish high-confidence divergence (large $\delta_i$, small
$\sigma_i$) from uncertain inference (large $\delta_i$, large $\sigma_i$).
Section~\ref{sec:inference} elaborates on the treatment of ill-posedness;
Section~\ref{sec:evaluation} quantifies the empirical regime in which IOR's
posterior uncertainty is small enough to constitute useful regulatory evidence.
